\documentclass[11pt,a4paper]{article}
\usepackage[T1]{fontenc}
\usepackage{geometry}
\usepackage{amsmath}
\usepackage{amssymb}
\usepackage{booktabs}
\usepackage{graphicx}
\usepackage{xcolor}
\usepackage{array}
\usepackage{mathptmx}
\usepackage{listings}
\usepackage{hyperref}

\geometry{margin=1in}

\definecolor{codebg}{gray}{0.95}
\lstset{
  basicstyle=\ttfamily\small,
  backgroundcolor=\color{codebg},
  breaklines=true,
  frame=single,
  columns=fullflexible,
  showstringspaces=false
}

\title{Forcing FP Exceptions in HeCBench's \texttt{fluidSim-cuda}\\
\large A Deliberately Unstable LBM Variant, Observed Under Nixnan}
\author{SC26 Tutorial / Nixnan Experiments}
\date{}

\begin{document}
\maketitle

\section{Executive Summary}

\texttt{hec\_fluidSim/} is a local copy of HeCBench's \texttt{fluidSim-cuda}
benchmark (a D2Q9 Lattice-Boltzmann Method fluid solver), kept as this
tutorial's own test area since \texttt{external/HeCBench} is a submodule we
cannot push changes to. Run as shipped, this benchmark never produces a
single FP exception --- its default relaxation parameter is safely inside
the scheme's stability bound. This report documents:

\begin{enumerate}
  \item \textbf{How a genuinely exception-producing variant was built}
    (\texttt{main\_overflow.cu}): pushing the BGK relaxation parameter
    \texttt{omega} far outside its $0 < \omega < 2$ stability bound turns
    the per-iteration collision update from a damping recurrence into an
    amplifying one, overflowing double's range within about 40 iterations
    and driving the local density through exactly zero, which the
    kernel's own undivided-by-zero-guarded \texttt{u.x = (\ldots)/rho}
    then turns into a second, independent exception source
    (Section~\ref{sec:mechanism}).
  \item \textbf{The result, confirmed two ways}: a standalone host-side
    scan (580,644 of 589,824 output cells are NaN) and an independent
    unsampled nixnan run reporting the same failure at the SASS
    instruction level --- 1,249 distinct FP64 NaN sites (16.7~million
    repeats) in the kernel's own \texttt{DFMA}/\texttt{DMUL}/\texttt{DADD}
    instructions (Section~\ref{sec:results}).
  \item \textbf{Two binary-level false positives}, both real, but with
    different, distinguishable root causes (Section~\ref{sec:falsepos}):
    one is a second confirmed instance of an already-documented
    \texttt{lu\_solve\_gpu} finding (a compiler register-move idiom
    misread as genuine \texttt{f32} arithmetic); the other is a new,
    only partially diagnosed anomaly in nixnan's own type-reporting
    machinery, where a definitively \texttt{f64}-classified instruction
    (\texttt{DFMA}, confirmed directly against nixnan's instruction
    table) is nonetheless \emph{labeled} \texttt{f16} in the live error
    report.
\end{enumerate}

\section{Provenance and Algorithm}

\texttt{fluidSim-cuda} is a D2Q9 Lattice-Boltzmann Method (LBM) fluid
simulation on a fixed $256\times256$ 2D grid, copied from
\href{https://github.com/zjin-lcf/HeCBench}{zjin-lcf/HeCBench},
path \texttt{src/fluidSim-cuda}, commit \texttt{2b73f4d} (see
\texttt{hec\_fluidSim/\allowbreak README.md} for full attribution and licensing ---
the benchmark itself traces to a 2009 AMD APP SDK sample). Each grid cell
holds a particle-distribution function over the 9 D2Q9 discrete
velocities. One CUDA kernel, \texttt{lbm} (\texttt{kernels.cu}), fuses the
two classic LBM steps every iteration:

\begin{itemize}
  \item \textbf{Collision}: for fluid (non-boundary) cells, computes local
    density $\rho$ and velocity $u$ from the current distributions, then
    relaxes each distribution toward its local equilibrium value
    (\texttt{ced()}, implementing
    $f_{\mathrm{eq}} = \rho\, w\,(1 + 3(e{\cdot}u) + 4.5(e{\cdot}u)^2 -
    1.5|u|^2)$) via
    \[
    f_{\mathrm{new}} = (1-\omega)\,f_{\mathrm{old}} + \omega\,f_{\mathrm{eq}}.
    \]
  \item \textbf{Streaming}: pushes each cell's post-collision values to
    its 8 neighbors along the corresponding D2Q9 direction, with the
    domain's boundary ring of cells (\texttt{h\_type}) handled via
    bounce-back instead.
\end{itemize}

The driver, \texttt{fluidSim()}, launches \texttt{lbm} once per requested
iteration, ping-ponging between two distribution buffers each step.
Upstream's own \texttt{main.cu} uses $\omega = 1.2$ and a small initial
velocity ($u_0 = (0.01, 0.01)$) --- both safely inside the scheme's stable
regime, which is exactly why the unmodified benchmark never produces an
FP exception under nixnan, regardless of iteration count.

\section{How the NaN-Producing Variant Was Made}
\label{sec:mechanism}

\subsection{The stability bound, and how to break it on purpose}

The collision update $f_{\mathrm{new}} = (1-\omega) f_{\mathrm{old}} +
\omega f_{\mathrm{eq}}$ is a textbook damped recurrence: any deviation
$\delta = f_{\mathrm{old}} - f_{\mathrm{eq}}$ between a cell's actual and
equilibrium distributions carries forward as $(1-\omega)\,\delta$ each
step. For $0 < \omega < 2$, $|1-\omega| < 1$ and deviations shrink every
iteration --- this is why the scheme is stable at upstream's $\omega =
1.2$. Push $\omega$ far outside that range and $|1-\omega|$ becomes a
large multiplier instead: deviations grow roughly $|1-\omega|\times$ per
iteration.

\texttt{main\_overflow.cu} (a copy of \texttt{main.cu}, with the
\texttt{\#ifdef VERIFY} reference-comparison path removed --- see
Section~\ref{sec:results} for why a GPU/CPU comparison stops being
meaningful here) sets
\[
\omega = 10^{8}
\]
--- about $8\times10^{7}$ past the stability bound. The simulation's
initial state is constructed to already sit \emph{exactly} at equilibrium
(the same \texttt{computefEq()} call fills both the initial distributions
and, implicitly, their own equilibrium target), so the very first
collision is a no-op regardless of $\omega$; deviations only start
appearing once streaming and boundary bounce-back mix values between
neighboring cells, from iteration~2 onward. Once they do, at
$|1-\omega|\approx10^{8}$ per iteration, an initial deviation on the
order of $10^{-3}$ reaches double's $\sim\!1.8\times10^{308}$ ceiling in
\[
\log_{10^8}\!\left(\frac{10^{308}}{10^{-3}}\right) \approx \frac{311}{8}
\approx 39
\]
iterations. \texttt{MakeOverflow}'s \texttt{run} target uses 60 for
margin.

\subsection{A second exception source: division by (an oscillating) density}

\texttt{kernels.cu}'s \texttt{lbm} kernel computes velocity from density
with no zero-guard:
\begin{lstlisting}[language=C++]
u.x = (dot(f1234, x1234) + dot(f5678, x5678)) / rho;
u.y = (dot(f1234, y1234) + dot(f5678, y5678)) / rho;
\end{lstlisting}
Because $(1-\omega)$ is a large \emph{negative} number once $\omega > 1$,
the collision update flips sign every iteration as it diverges --- so the
locally-summed density $\rho$ (itself a sum of the nine distributions)
does not just grow in magnitude, it oscillates in sign while doing so,
and is expected to cross exactly $0.0$ at some cell, some iteration,
before overflow finishes the job. That turns \texttt{u.x}/\texttt{u.y}
into a direct $x/0$ division --- a second, independent route to
Infinity/NaN, on top of plain multiplicative overflow.

\subsection{Build}

\begin{lstlisting}[language=bash]
# MakeOverflow (new here; see README.md) -- no -DVERIFY
make -f MakeOverflow run   # ./main_overflow 60
\end{lstlisting}

\section{Results}
\label{sec:results}

\subsection{Standalone sanity check (independent of nixnan)}

\texttt{main\_overflow.cu} scans its own output for non-finite cells
before nixnan ever enters the picture:
\begin{center}
\begin{tabular}{lr}
\toprule
Output cells & Count \\
\midrule
Finite & 9{,}180 \\
NaN & 580{,}644 \\
Infinity & 0 \\
\midrule
Total & 589{,}824 \\
\bottomrule
\end{tabular}
\end{center}
98.4\% of all output cells are NaN after 60 iterations at $\omega=10^8$ ---
confirming the mechanism works, independent of any instrumentation. (No
raw Infinity survives to the \emph{final} output buffer; see
Section~\ref{sec:falsepos-intro} below for why nixnan's own live report,
taken \emph{during} the run rather than at the end, does show Infinity.)
This is also why \texttt{main\_overflow.cu} does not use \texttt{-DVERIFY}:
once both the GPU and a CPU \texttt{reference()} implementation are
diverging chaotically, floating-point non-associativity makes them
diverge in different, equally meaningless ways --- a GPU/CPU numerical
match stops being a meaningful question once the scheme itself is
unstable.

\subsection{Nixnan (unsampled, default exception checking)}

\begin{lstlisting}[language=bash]
LD_PRELOAD=../external/nixnan/nixnan.so ./main_overflow 60 \
  > Log/fluidSim-overflow-60iter-nixnan-default-exceptions.log 2>&1
\end{lstlisting}

\begin{center}
\begin{tabular}{lrrrr}
\toprule
Format & NaN (repeats) & Infinity (repeats) & $-$Infinity (repeats) & Subnormal (repeats) \\
\midrule
FP16 & 17 (402{,}319) & 8 (1{,}067) & 0 (0) & 0 (0) \\
FP32 & 6 (268{,}218) & 0 (0) & 0 (0) & 0 (0) \\
FP64 & 1{,}249 (16{,}701{,}143) & 831 (433{,}209) & 824 (444{,}362) & 26 (28{,}314) \\
\bottomrule
\end{tabular}
\end{center}
Distinct-site counts are how many unique (instruction, operand) pairs
first triggered each exception type; repeat counts are total occurrences
across the whole 60-iteration run. The FP64 row is the genuine result:
1,249 distinct sites and 16.7~million total NaN occurrences in the
kernel's own \texttt{DFMA}/\texttt{DMUL}/\texttt{DADD} instructions ---
exactly \texttt{lbm}'s own collision/streaming arithmetic, all
double-precision, as expected. The 26 distinct Subnormal sites are
consistent with values genuinely passing through the full extreme-
magnitude spectrum (subnormal $\to$ normal $\to$ overflow $\to$ NaN) as
the amplification runs away. The FP16 and FP32 rows are false positives,
covered next.

\section{A Binary-Level False Positive, Revisited}
\label{sec:falsepos}
\label{sec:falsepos-intro}

\texttt{lbm} is written entirely in \texttt{double}; every genuine
computation in it is FP64. The nonzero FP16/FP32 rows above are therefore
both false positives in the sense already documented for
\texttt{lu\_solve\_gpu}
(\texttt{lu\_solve\_\allowbreak gpu/\allowbreak
lu\_solve\_gpu\_\allowbreak report.tex}, ``An FP32 Artifact of
Binary Instrumentation''): \emph{a nonzero count in a format the source
never uses does not necessarily mean that format's arithmetic ran.} But
the two rows have different, and only one of them fully understood,
root causes.

\subsection{The FP32 case: the same register-move idiom, confirmed again}

Every FP32-classified error line in this run is one of exactly two
instruction texts:
\begin{lstlisting}
FFMA R6, RZ, R5, R15
FFMA R6, RZ, R5, R33
\end{lstlisting}
This is precisely \texttt{lu\_solve\_gpu}'s documented pattern:
\texttt{ptxas}'s ``multiply-by-the-zero-register'' idiom
($0\times R5 + R15 = R15$), a plain register \emph{move}, here used to
shuffle one 32-bit half of a 64-bit double across the register file ---
not a genuine 32-bit computation. Because nixnan classifies by SASS
opcode (\texttt{FFMA} $\to$ \texttt{f32}, unconditionally, per
\texttt{src/instructions.h}), it correctly identifies the instruction as
\texttt{f32}-typed and correctly notices the operand's bit pattern
decodes as NaN --- it is simply never told that this particular
\texttt{FFMA} is standing in for a double-register move, not real
\texttt{float} arithmetic. This is a second, independent confirmation of
the \texttt{lu\_solve\_gpu} finding, in an entirely different codebase.

\subsection{The FP16 case: a distinct, only partially diagnosed anomaly}

Every FP16-classified error line, by contrast, is one of exactly nine
instruction texts, all of the same shape:
\begin{lstlisting}
DFMA R8,  R14, -1.5, R8
DFMA R10, R14, -1.5, R10
DFMA R12, R14, -1.5, R12
DFMA R14, R14, -1.5, R6
DFMA R36, R14, -1.5, R36
DFMA R38, R14, -1.5, R38
DFMA R40, R14, -1.5, R40
DFMA R42, R14, -1.5, R42
DFMA R50, R14, -1.5, R50
\end{lstlisting}
This is \emph{not} the same phenomenon. \texttt{DFMA} is not an idiom
here: the constant \texttt{-1.5} is exactly the $-1.5|u|^2$ coefficient
in \texttt{ced()}'s equilibrium formula (Section~\ref{sec:mechanism}), and
\texttt{R14} plausibly holds $|u|^2$ --- this is genuine, textbook
double-precision arithmetic computing part of the equilibrium
distribution, not a register-shuffle trick. Directly decoding nixnan's
own instruction table (\texttt{src/instructions.h}, extracted from its
byte-array encoding) confirms this: the \texttt{DFMA} entry classifies
\emph{all four} operand positions as \texttt{f64}, unconditionally, with
no \texttt{f16}-typed \texttt{DFMA} pattern anywhere in the table:
\begin{lstlisting}
"(DFMA(?:\.(RP|RM|RZ))?)": {
  "registers": [
    {"type": "f64", "count": 1}, {"type": "f64", "count": 1},
    {"type": "f64", "count": 1}, {"type": "f64", "count": 1}
  ],
  "description": "Double-precision fused multiply-add with optional FTZ and saturation"
}
\end{lstlisting}
So the table itself never says \texttt{DFMA} is \texttt{f16}. Reading
\texttt{nixnan.cu}'s error-reporting code (around the
\texttt{recv\_thread\_fun} loop) shows the \texttt{type} string printed
per report is \emph{not} looked up from the table at report time; it
comes from
\begin{lstlisting}[language=C++]
type = type_to_string.at(recorder->get_type(id, ei->operand()));
\end{lstlisting}
--- a runtime lookup keyed by an internal instruction \texttt{id} and the
operand slot, populated once when the instruction was first instrumented.
That the table unambiguously says \texttt{f64} while this lookup
sometimes returns \texttt{f16} for the same instruction text points to
the \texttt{id}$\to$type mapping itself, not the classification table, as
the likely site of the bug --- e.g.\ an \texttt{id} collision between this
\texttt{DFMA} and an unrelated, genuinely \texttt{f16}-typed instruction
elsewhere in the same kernel, so that the wrong stored type is retrieved
for this operand. \textbf{This report does not go further than that}: the
detection itself is not in question (a real NaN genuinely appears at this
operand, at this instruction, in this many-times-repeated pattern) ---
only the \texttt{f16} \emph{label} attached to it is wrong, and exactly why
is not fully diagnosed here. Worth raising with nixnan's maintainers as a
reproducible case: this exact 60-iteration run, this exact instruction
list.

\section{Conclusion}

\begin{enumerate}
  \item \texttt{fluidSim-cuda}'s LBM scheme is stable, and exception-free
    under nixnan, entirely because of its default relaxation parameter
    choice --- not because the kernel is somehow immune to FP exceptions.
    Pushing $\omega$ outside its $0 < \omega < 2$ stability bound reliably
    produces both overflow (from the amplifying collision recurrence) and
    division-by-(oscillating-through-)zero (from the kernel's unguarded
    \texttt{u.x = (\ldots)/rho}), confirmed both by a standalone host-side
    scan and by nixnan's own instruction-level report.
  \item Binary-level FP exception detection can still be trusted at the
    level that matters most --- \emph{a real exception really occurred,
    at this instruction} --- even when its \emph{format label} is wrong.
    Two false-positive format labels were found in this one run, with two
    different causes: one already-documented and well-understood (a
    compiler register-move idiom, \texttt{lu\_solve\_gpu}'s finding,
    reproduced here independently); one new and only partially diagnosed
    (a likely \texttt{id}-collision in nixnan's own type-recording
    machinery, isolated here to a specific, reproducible instruction
    pattern but not fully root-caused).
\end{enumerate}

\section*{Reproduction}

\begin{lstlisting}[language=bash]
cd hec_fluidSim
make -f MakeOverflow run
LD_PRELOAD=../external/nixnan/nixnan.so ./main_overflow 60 \
  > Log/fluidSim-overflow-60iter-nixnan-default-exceptions.log 2>&1
\end{lstlisting}
Requires a built \texttt{../external/nixnan/nixnan.so}. See
\texttt{README.md} in this directory for provenance, licensing, and a
description of every file (including \texttt{MakeShort} and
\texttt{run\_reduced.sh}, unrelated to this report's exception-forcing
variant).

\end{document}
